\documentclass[12pt,a4paper]{article}

\usepackage{amsmath,amssymb,bm}
\usepackage[utf8]{inputenc}
\usepackage[T1]{fontenc}
\usepackage{geometry}
\usepackage{booktabs}
\usepackage{enumitem}
\usepackage{hyperref}
\usepackage{xcolor}
\usepackage{tabularx}
\usepackage{array}
\usepackage{parskip}
\usepackage[most]{tcolorbox}
\usepackage{setspace}

% ---- pseudocode box identical in spirit to the FOIL-algorithm image ----
\newtcolorbox{algobox}[1][]{%
  colback=white, colframe=black,
  boxrule=0.6pt, arc=0pt,
  left=6pt, right=6pt, top=4pt, bottom=4pt,
  before upper={\setstretch{1.25}\small},
  #1}

% Convenience macros for pseudocode typography
\newcommand{\kw}[1]{\textbf{#1}}          % keyword
\newcommand{\fn}[1]{\textsc{#1}}          % function name (small caps)
\newcommand{\var}[1]{\textit{#1}}         % variable name
\newcommand{\cm}[1]{\textcolor{gray}{#1}} % comment
\newcommand{\ind}{\hspace*{1.4em}}        % one indent level
\newcommand{\indd}{\hspace*{2.8em}}       % two indent levels
\newcommand{\inddd}{\hspace*{4.2em}}      % three indent levels
\newcommand{\assign}{\ensuremath{\leftarrow}}%

\geometry{margin=2.5cm}

\hypersetup{
    colorlinks=true,
    linkcolor=blue,
    citecolor=blue,
    urlcolor=blue
}

\definecolor{mplanned}{rgb}{0.0,0.45,0.70}
\definecolor{murgent}{rgb}{0.80,0.10,0.10}
\definecolor{mminor}{rgb}{0.10,0.50,0.10}

\newcommand{\planned}[1]{{\color{mplanned}#1}}
\newcommand{\urgent}[1]{{\color{murgent}#1}}
\newcommand{\minor}[1]{{\color{mminor}#1}}

\title{\textbf{Revision Plan}\\[4pt]
\large Manuscript SAMO-D-26-00346\\
\textit{``Frequency maximization of planar structures using\\quasi-static approximation in topology optimization''}}

\author{Piotr Tauzowski \and Andrzej Szczepańczyk \and Bartłomiej Błachowski}

\date{May 2026}

\begin{document}

\maketitle
\tableofcontents
\newpage

% ============================================================
\section{Overview and Prioritisation}
% ============================================================

Three reviewers provided detailed comments on the manuscript.
Reviewer~1 raises fundamental concerns about novelty framing and requests
additional robustness experiments. The Reviewer~3 PDF report identifies
important issues about the source of computational efficiency and flags
several specific items (missing references, unclear figures, grey regions).
Reviewer~4 asks four focused technical questions about the formulation.

Table~\ref{tab:priority} summarises all requested changes, grouped by
priority.  Items marked \urgent{urgent} must be resolved before resubmission
because they concern scientific correctness or the main claim of the paper.
Items marked \planned{planned} are required for acceptance. Items marked
\minor{minor} are easily addressed and can be completed in one pass.

\begin{table}[htbp]
\centering
\caption{Priority map of required revisions (R = Reviewer, \# = comment number).}
\label{tab:priority}
\renewcommand\arraystretch{1.15}
\begin{tabularx}{\textwidth}{@{}p{0.7cm}p{1.6cm}p{8.5cm}p{2.5cm}@{}}
\toprule
ID & Source & Issue & Priority \\
\midrule
A0 & Internal audit & Align manuscript equations with implemented frozen-load definition
(baseline model, $\omega$ vs.\ $\omega^2$, and sensitivity path) & \urgent{Urgent} \\
A1 & R1 Major & Reframe novelty vs.\ Yuksel--Yilmaz & \urgent{Urgent} \\
A2 & R1 Major & Table 1: break down runtime (setup, per-iteration, iterations) & \urgent{Urgent} \\
A3 & R1 Major & Explain 1083/1723-iteration discrepancy vs.\ Yuksel's own paper & \urgent{Urgent} \\
A4 & R1 Major & Approximate-error analysis; N-iteration eigenpair refresh study & \urgent{Urgent} \\
A5 & R1 Major & Mode clustering / coalescence behaviour; a~posteriori mode check & \urgent{Urgent} \\
A6 & R3 Gen.  & Clarify that speedup comes from convergence rate, not per-iter cost & \urgent{Urgent} \\
\midrule
B1 & R3 \#8  & Explain two new modes at $\alpha=1$: are they localized / spurious? & \planned{Planned} \\
B2 & R3 \#9  & High-index mode tracking for $\Phi_2$ — discuss artificial modes & \planned{Planned} \\
B3 & R3 \#10 & Apply B1 and B2 discussion to building example (Section~3.3) & \planned{Planned} \\
B4 & R3 \#11 & RAMP interpolation: consider for final eigenvalue evaluation & \planned{Planned} \\
B5 & R3 \#12 & Grey regions in Figs.~9 and~10 — comment on convergence / filter & \planned{Planned} \\
B6 & R4 \#3  & Clarify why frozen-load approach avoids spurious low-density modes & \planned{Planned} \\
\midrule
C1 & R3 \#1  & Add application references in Introduction paragraph~1 & \minor{Minor} \\
C2 & R3 \#2  & Add canonical refs for SIMP, BESO, LSM, MMC at first mention & \minor{Minor} \\
C3 & R3 \#3  & Add Rayleigh's principle reference + brief introductory paragraph & \minor{Minor} \\
C4 & R3 \#4  & Add MMA reference in Section~2.1 & \minor{Minor} \\
C5 & R3 \#5  & State hardware platform details for Table~1 & \minor{Minor} \\
C6 & R3 \#6  & Add filter equation + reference in Section~2.1 (or cite Sec.~2.4) & \minor{Minor} \\
C7 & R3 \#7  & Show sensitivity expression for problem~(4) explicitly in Section~2.2 & \minor{Minor} \\
C8 & R3 \#13 & Clarify what supplementary material contains & \minor{Minor} \\
C9 & R4 \#1  & Clarify load-sensitivity treatment (lines 798--799) & \minor{Minor} \\
C10& R4 \#2  & Discuss potential numerical difficulties from design-dependent load & \minor{Minor} \\
C11& R4 \#4  & Comment on symmetric vs.\ asymmetric domain applicability & \minor{Minor} \\
\bottomrule
\end{tabularx}
\end{table}

\newpage
% ============================================================
\section{Urgent Revisions (A-group)}
% ============================================================

\subsection{A0 — Align the manuscript with the implemented frozen-load definition}

\textbf{Internal audit finding.}  Before addressing the reviewer comments,
the authors must resolve a code--manuscript mismatch in the definition of the
proposed load.  The manuscript describes
$\mathbf{f}=(\omega_j^{(0)})^2\mathbf{M}(\mathbf{x})\boldsymbol{\Phi}_j$
with $(\omega_j^{(0)},\boldsymbol{\Phi}_j)$ computed on the initial
uniform-density design.  The submitted example configurations, however, use
\texttt{semi\_harmonic} loads with
\texttt{semi\_harmonic\_baseline="solid"}; in the current MATLAB and Python
implementations this assembles
\[
  \mathbf{f}_{\mathrm{semi}}(\mathbf{x})
  = \rho_{\mathrm{nodal}}(\mathbf{x})\,
    \bigl(\omega_j^{(0)}\mathbf{M}_0\boldsymbol{\Phi}_j\bigr),
\]
i.e. the baseline matrices are solid/reference matrices and the stored
modal-load vector uses $\omega_j^{(0)}$, not $(\omega_j^{(0)})^2$.
This is not a harmless documentation detail for the two-mode examples,
because the relative weighting of mode~1 and mode~2 changes if $\omega$ is
replaced by $\omega^2$.

\textbf{Required correction.}
\begin{itemize}[leftmargin=1.6em]
  \item Choose one definition and make the paper, JSON examples, MATLAB code,
        Python code, and response letter consistent.  If the intended method
        is the physical inertial-load approximation, change the
        \texttt{semi\_harmonic} base vector to
        $(\omega_j^{(0)})^2\mathbf{M}_0\boldsymbol{\Phi}_j$ and rerun the
        reported topologies, MAC tables, and performance table.  If the
        current implementation is retained, revise the manuscript equations
        and claims to call it a normalized/reference modal load rather than
        the physical inertial force.
  \item Resolve the baseline-model statement.  The current examples compute
        $\boldsymbol{\Phi}_j$ on the solid/reference design, not the uniform
        $x_e=V_f$ design.  Either rerun with
        \texttt{semi\_harmonic\_baseline="initial"} or revise the manuscript
        consistently to state that the reference modes are taken from the
        fully solid domain.
  \item Keep the sensitivity statement tied to the actual load type:
        \texttt{semi\_harmonic} intentionally zeroes
        $\partial\mathbf{f}/\partial x_e$; \texttt{harmonic} currently adds a
        load-sensitivity contribution.  The \texttt{harmonic} path should not
        be cited as evidence for the frozen-load sensitivity unless an option
        is added to disable that term.
\end{itemize}

\subsection{A1 — Reframe novelty relative to Yuksel--Yilmaz}

\textbf{Reviewer~1 concern.}  The reviewer argues that the proposed method is
merely a simplification of the Yuksel--Yilmaz iterative scheme (replacing the
self-consistent eigenvector update with a fixed estimate), rather than a
fundamentally different approach. The current abstract and introduction may
overstate the novelty.

\textbf{What needs to change.}
\begin{itemize}[leftmargin=1.6em]
  \item Revise the abstract and the end of the Related Work section
        (Section~2.4, last two paragraphs) to state the relationship
        \emph{explicitly}: after resolving A0, the proposed method should be
        described as a further simplification of the Yuksel--Yilmaz static
        approximation framework in which the reference mode-shape estimate is
        computed exactly \emph{once} (on the declared reference design) and
        then frozen permanently, whereas Yuksel--Yilmaz update the estimate at
        every iteration through a design-dependent inertial load.
  \item Reframe the novelty as (i)~the frozen-mode variant eliminates the
        load-design coupling entirely, reducing each iteration to an uncoupled
        compliance step with no additional eigenproblem machinery inside the
        loop; (ii)~the multi-mode weighting parameter~$\alpha$ and its
        MAC-verified monotone gain rule constitute a distinct, demonstrated
        contribution not present in Yuksel--Yilmaz.
  \item Add a sentence acknowledging that the method sacrifices the
        self-consistency of the mode estimate and thereby introduces an
        accuracy ceiling, which is the known trade-off of the simplification.
\end{itemize}

\textbf{Source code context.}  The submitted examples use the configured
\texttt{semi\_harmonic} load path in
\texttt{analysis/ourApproach/Matlab/topopt\_freq.m} and
\texttt{analysis/ourApproach/Python/topopt\_freq.py}, not the legacy
\texttt{harmonic} fallback.  In that path the reference eigenpairs are cached
once, the mode shapes are never updated during the loop, and
\texttt{\_load\_sensitivity}/\texttt{localLoadSensitivity} returns no
load-sensitivity contribution for \texttt{semi\_harmonic}.  The revised text
should therefore describe the active \texttt{semi\_harmonic}
frozen-reference path, subject to the load-definition correction in A0.

% -----------------------------------------------------------
\subsection{A2 — Runtime breakdown in Table~1}

\textbf{Reviewer~1 and Reviewer~3 concern.}  The table reports only total
runtime, total iterations, and peak RAM. Reviewer~1 requests a breakdown
separating (i)~one-time initialization cost (eigensolve at Step~1),
(ii)~per-iteration cost, and (iii)~total iteration count so that the source
of the speedup is identifiable. Reviewer~3 makes the same point more sharply:
the proposed method's per-iteration time is \emph{higher} than that of the
Yuksel code at the two largest meshes (0.08 vs.\ 0.05~s/iter at
$320{\times}40$; 0.12 vs.\ 0.08~s/iter at $400{\times}50$), so the net
speedup comes entirely from fewer iterations.

\textbf{What needs to change.}
\begin{itemize}[leftmargin=1.6em]
  \item First rerun timing benchmarks with diagnostic frequency-history
        eigensolves disabled.  The current performance driver disables live
        visualization and image saving, but the base JSON has
        \texttt{postprocessing.save\_frequency\_iterations=true}.  In the
        proposed MATLAB solver this option explicitly forces per-iteration
        eigensolves for plotting and makes timing comparisons unfair.
        Generate frequency-history plots in a separate diagnostic run.
  \item Add a column ``Run time/iter (s)'' to Table~1 if not already
        present, \emph{and} an additional column ``Init.\ time (s)''
        reporting the one-time eigensolve cost at initialization.  (The
        current table already has the per-iteration column; make sure the
        init time is separated.)
  \item In the text discussing Table~1 (Section~4.1 / 4.2), add explicit
        statements such as: ``The proposed method has a higher per-iteration
        cost than the Yuksel code at fine meshes (\ldots~s vs.\ \ldots~s),
        but converges in substantially fewer iterations (\ldots\ vs.\
        \ldots), yielding a lower total runtime. The speedup is therefore
        a consequence of improved convergence behaviour, not of reduced
        per-iteration cost.''
  \item Provide a brief mechanistic explanation: the frozen, consistent
        descent direction from the first iteration prevents the oscillatory
        behaviour that appears to inflate the Yuksel iteration count at fine
        meshes.
\end{itemize}

\textbf{Source code context.}  The per-iteration timing and total iteration
count are already tracked (\texttt{tIter}/\texttt{nIter} in MATLAB and
\texttt{t\_iter}/\texttt{n\_iter} in Python).  Initialization timing must
cover the actual active path: \texttt{semi\_harmonic} baseline eigensolves,
\texttt{harmonic} first-update eigensolves, and the legacy fallback.  The
MATLAB JSON wrapper currently returns only \texttt{x}, \texttt{omega},
\texttt{tIter}, \texttt{nIter}, and memory, so the diagnostics interface must
also be extended if Table~1 is generated from MATLAB.

% -----------------------------------------------------------
\subsection{A3 — Yuksel iteration-count discrepancy}

\textbf{Reviewer~1 concern.}  The paper reports 1723 iterations for the
Yuksel code at $320{\times}40$ and 1083 at $400{\times}50$, whereas
Yuksel--Yilmaz's own paper reports approximately 180--200 iterations for the
same $320{\times}40$ mesh. Even accounting for different hardware and
convergence criteria, a $8{\times}$--$10{\times}$ difference in iteration
count is not explained.

\textbf{What needs to change.}
\begin{itemize}[leftmargin=1.6em]
  \item Add a paragraph (or extend the existing algorithmic comparison in
        Section~4.1) that explains the discrepancy. Possible reasons:
        (a)~the convergence criterion used by the authors (\,$\max|\Delta
        x_e|<10^{-3}$) may be stricter or different from what Yuksel--Yilmaz
        use; (b)~the iteration count in this repository is the sum of
        Stage~1 and Stage~2 iterations, whereas the reference paper may
        report only the second stage; (c)~Heaviside post-processing,
        frequency-history diagnostics, filter settings, or convergence
        criteria may differ from the reference code; (d)~the Yuksel code as
        re-implemented by the authors may use implementation choices not
        present in the original.
  \item State which implementation was used (the authors' own MATLAB
        reimplementation in \texttt{analysis/YukselApproach/Matlab}, not the
        original released code unless that is obtained and run) and clarify
        whether reported iteration counts include Stage~1, Stage~2, or both.
  \item If the discrepancy cannot be fully resolved, state it as an open
        question and note that the comparison is therefore between
        implementations, not between the methods as originally published.
\end{itemize}

% -----------------------------------------------------------
\subsection{A4 — Approximation error analysis and eigenpair refresh study}

\textbf{Reviewer~1 concern.}  The paper does not quantify the approximation
error introduced by freezing the eigenvector. The reviewer suggests:
(a)~providing quantitative bounds on the error; and (b)~a numerical
sensitivity study showing how the final frequency and topology depend on the
initial eigensolve, for example by recomputing the eigenpair every $N$
iterations.

\textbf{What needs to change.}
\begin{itemize}[leftmargin=1.6em]
  \item \textbf{Theoretical framing.}  Add a short paragraph (possibly in
        Section~2.2 after Eq.~4) discussing Rayleigh's principle as a bound:
        since $\omega_1^2 \leq R(\boldsymbol{\Phi}_1)$ with equality only
        when $\boldsymbol{\Phi}_1$ equals the true mode, the gap
        $R(\boldsymbol{\Phi}_1) - \omega_1^2$ measures the error in the
        mode shape.  As the MAC value of the tracked mode is~0.9998 for the
        simply-supported beam, the mode-shape error is demonstrably small for
        this case.
  \item \textbf{Numerical study: eigenpair refresh.}  Run the simply-supported
        beam benchmark at $400{\times}50$ with the eigenpair refreshed every
        $N \in \{1, 10, 50, \infty\}$ iterations. Report final
        $\omega_1$, runtime, and iteration count in a new small table. This
        quantifies the accuracy--cost tradeoff and directly addresses the
        reviewer's request.
        \emph{Implementation caveat:} the current \texttt{update\_after}
        schedule applies only to \texttt{harmonic} loads, whereas the reported
        proposed examples use \texttt{semi\_harmonic}.  The \texttt{harmonic}
        branch also includes a load-sensitivity term.  Therefore this study
        needs a controlled implementation path: recompute the reference
        eigenpair every $N$ iterations while keeping the same load formula and
        the same zero-load-sensitivity approximation as the proposed method.
        Do not present \texttt{update\_after=1} as the Yuksel method; it is a
        true-eigenpair refresh variant, not Yuksel--Yilmaz's static
        self-consistent update.
  \item Place this study in a new subsection (e.g.\ Section~4.3
        ``Sensitivity to eigenpair update frequency'') or as an appendix.
\end{itemize}

% -----------------------------------------------------------
\subsection{A5 — Mode clustering, coalescence, and a~posteriori mode check}

\textbf{Reviewer~1 concern.}  The paper discusses mode-ordering shifts and
uses MAC tracking but does not explain how the method behaves when eigenvalues
are clustered or when the target mode changes shape during the optimization.
The reviewer also requests (a)~an a~posteriori check that the targeted mode
remains the structure's lowest mode after convergence, and (b)~a clear
statement that the method targets a specific mode shape rather than the
fundamental frequency.

\textbf{What needs to change.}
\begin{itemize}[leftmargin=1.6em]
  \item Add a paragraph in Section~2.5 (Mode-shape tracking via MAC) or in
        the Conclusions explicitly stating that the quasi-static load
        $\mathbf{f} \propto \mathbf{M}(\mathbf{x})\boldsymbol{\Phi}_1$
        targets the mode that \emph{retains the shape} $\boldsymbol{\Phi}_1$,
        not necessarily the lowest mode of the final topology.  Mode ordering
        can shift (as observed at $\alpha=1$ in both beam and building
        examples) and this is expected and acceptable behavior.
  \item Document the a~posteriori mode check: after convergence, the
        algorithm already solves for the lowest three modes (Step~5 in
        Algorithm~1). Report the MAC between each of these modes and
        $\boldsymbol{\Phi}_1$, and identify the tracked mode. This is already
        done in the paper (Tables~2--5) but should be explicitly connected to
        the a~posteriori check the reviewer requests.
  \item Discuss mode coalescence briefly: the frozen-mode approach does not
        explicitly enforce coalescence (unlike the Olhoff bound formulation)
        and does not suffer from the differentiability issues of repeated
        eigenvalues precisely because no eigenvalue is differentiated inside
        the loop.  However, near-coalescence at the final design (as at
        $\alpha=0.75$ for the clamped beam, $\omega_1=347.7$ vs.\
        $\omega_2=360.4$~rad/s) may cause the tracked mode to interchange
        with a neighbor in a final eigensolve; MAC tracking detects this.
\end{itemize}

% -----------------------------------------------------------
\subsection{A6 — Reframe the efficiency claim}

\textbf{Reviewer~3 concern.}  The paper presents computational efficiency as
a primary advantage but attributes it ambiguously. The reviewer notes that the
per-iteration cost of the proposed method is \emph{higher} than for the Yuksel
code at the fine meshes, so the efficiency gain is driven by faster
convergence (fewer iterations), not by reduced per-iteration cost.

\textbf{What needs to change.}
\begin{itemize}[leftmargin=1.6em]
  \item In the abstract, replace ``requiring only a single eigensolve at
        initialization'' with language that clarifies \emph{two} sources of
        speedup: (i)~no eigensolve inside the iteration loop (removing
        Arnoldi/shift-invert overhead), and (ii)~consistently fewer
        iterations to convergence (stable, non-oscillating descent direction).
  \item In the discussion section (Section~4), add a short paragraph
        explicitly comparing per-iteration costs and iteration counts
        separately, and acknowledge that the per-iteration advantage over
        the Yuksel code is mesh-size-dependent and may reverse at fine
        meshes.
  \item Under which conditions the method is more efficient than
        Yuksel--Yilmaz should be stated explicitly: the proposed method has
        a net advantage whenever fewer iterations overcompensate the higher
        per-iteration cost.  In the current table this is true for the
        $240{\times}30$, $320{\times}40$, and $400{\times}50$ meshes, but not
        for $160{\times}20$ where Yuksel is slightly faster (4.7~s vs.\
        5.1~s).  This statement must be updated after rerunning timings with
        diagnostics disabled.  Whether this holds for 3D or very fine 2D
        meshes remains an open question to be noted in the Conclusions.
\end{itemize}

\newpage
% ============================================================
\section{Planned Revisions (B-group)}
% ============================================================

\subsection{B1 — Two new low-frequency modes at $\alpha=1$ (clamped beam)}

\textbf{Reviewer~3 comment~\#8.}  At $\alpha=1$, two structural modes appear
below the $\boldsymbol{\Phi}_1$-tracked mode (topo-mode~3, 368.2~rad/s).
Their MAC with any solid-domain shape is below 0.04. The reviewer asks
whether these are localized (spurious) modes associated with low-density
regions, and if so, whether this reveals a limitation regarding mode
switching.

\textbf{What needs to change.}
\begin{itemize}[leftmargin=1.6em]
  \item Inspect the mode shapes of topo-modes~1 and~2 for the $\alpha=1$
        clamped-beam topology (already available in the run data). If their
        spatial support is concentrated in low-density regions, they are
        spurious localized modes. If they span the structural skeleton, they
        are genuine new structural modes enabled by the mass redistribution.
  \item Add a clarifying sentence in Section~3.2: explain which of the two
        scenarios applies (localized vs.\ genuine), and state that MAC
        tracking correctly identifies them as not corresponding to any
        solid-domain shape, so the proposed tracking methodology is robust to
        this phenomenon.
  \item If they are spurious modes: note that they arise from using $d=p=3$
        for the mass interpolation in the proposed method (in contrast to
        the Olhoff implementation which uses $d=6$ for low-density elements,
        as stated in Section~2.1 of the paper). Reference the Reviewer~3
        comment~\#11 suggestion about RAMP (see B4 below).
\end{itemize}

% -----------------------------------------------------------
\subsection{B2 --- High-index mode tracking for \texorpdfstring{$\boldsymbol{\Phi}_2$}{Phi\_2}}

\textbf{Reviewer~3 comment~\#9.}  In Table~3 (clamped beam), the mode
best-matching $\boldsymbol{\Phi}_2$ appears at indices 18, 30, 40, 46, or 91
in the final eigenvalue ordering, suggesting a large number of low-frequency
modes that may obscure genuinely structural modes.

\textbf{What needs to change.}
\begin{itemize}[leftmargin=1.6em]
  \item Add a note in Section~3.2 explaining that the high mode indices
        for $\boldsymbol{\Phi}_2$ tracking are consistent with the presence
        of many spurious localized modes in the low-density regions (void
        elements with $x_e \approx 10^{-3}$ and $d=p=3$), which populate
        the lower end of the spectrum.
  \item Emphasise that MAC-based tracking bypasses the need to inspect each
        mode: only the mode with the highest MAC value (regardless of index)
        is reported, and its frequency is the physically relevant quantity.
  \item Consider adding mode shape visualisations for the MAC-best modes
        in supplementary material to make the claim concrete.
\end{itemize}

% -----------------------------------------------------------
\subsection{B3 — Same issues in building example (Section~3.3)}

The observations from B1 and B2 also manifest in the building example
(Table~5): the $\boldsymbol{\Phi}_2$-tracked mode reaches indices 33, 58, 69,
71, or 81. The same discussion should be added or cross-referenced in
Section~3.3.

% -----------------------------------------------------------
\subsection{B4 — RAMP interpolation for spurious-mode suppression}

\textbf{Reviewer~3 comment~\#11.}  The reviewer suggests considering the RAMP
interpolation scheme as an alternative to SIMP for the mass matrix, since RAMP
can be tuned to push low-density modes toward higher frequencies, reducing
spectral pollution.

\textbf{What needs to change.}
\begin{itemize}[leftmargin=1.6em]
  \item A full switch to RAMP in the optimization loop is beyond the current
        scope.  A post-processing eigensolve with an alternative mass
        interpolation may be useful, but it should not be described as a
        one-line or formula-free fix: the exact stiffness/mass interpolation
        pair must be taken from the cited RAMP literature and verified on the
        present plane-stress model before being reported.
  \item In the paper, add a sentence in Section~4.2 (proposed approach)
        acknowledging that SIMP with $d=p=3$ for the mass generates spurious
        localized modes in low-density regions (as also noted for the Olhoff
        implementation which uses $d=6$ for this reason). State that MAC
        tracking provides a robust mechanism to identify targeted modes
        independent of spectral pollution.  Discuss RAMP or $d>p$ as a
        plausible mitigation to be tested, not as an already-demonstrated
        explanation for the present results.
\end{itemize}

% -----------------------------------------------------------
\subsection{B5 — Grey regions in Figs.~9 and~10}

\textbf{Reviewer~3 comment~\#12.}  The optimized topologies for the clamped
beam (Fig.~9) and building (Fig.~10) contain visually grey regions with
intermediate densities.

\textbf{What needs to change.}
\begin{itemize}[leftmargin=1.6em]
  \item Add a sentence acknowledging the grey regions: the clamped-beam
        and building optimizations use the OC/MMA optimizer with sensitivity
        filtering but without Heaviside projection continuation
        (\texttt{use\_heaviside=False} in the JSON configs). The absence of
        projection allows intermediate densities to persist at convergence.
  \item State that adding a Heaviside continuation schedule (as in the
        Olhoff implementation) would sharpen the black-and-white
        transition, but at the cost of additional tuning parameters and
        more iterations. The current results are presented at the level of
        the sensitivity-filter baseline to maintain a fair algorithmic
        comparison with the Yuksel code.
  \item If time permits, run the clamped-beam example with Heaviside
        continuation and show the sharpened topology as a supplementary
        figure.  The Python wrapper already forwards the JSON
        \texttt{optimization.filter.heaviside} flag to \texttt{run\_cfg}, but
        the MATLAB wrapper currently appears to parse this flag only for the
        Olhoff path; if MATLAB is used for the paper results, forward the
        flag into the \texttt{ourApproach} \texttt{runCfg} before running this
        experiment.
\end{itemize}

% -----------------------------------------------------------
\subsection{B6 — Why the frozen-load approach avoids spurious localized modes}

\textbf{Reviewer~4 comment~\#3.}  The reviewer correctly observes that the
proposed approach may be less susceptible to the spurious-mode artefact than
standard dynamic formulations, and asks for clarification.

\textbf{What needs to change.}
\begin{itemize}[leftmargin=1.6em]
  \item Add a paragraph in Section~2.2 (or in Section~4.2): the
        quasi-static approach does not minimize an eigenvalue and therefore
        never solves the generalized eigenproblem inside the loop where
        spurious modes could corrupt sensitivity computation. Spurious
        localized modes \emph{do} appear in the final post-convergence
        eigensolve (Step~5), but at that point they only affect the
        eigenvalue spectrum and not the optimization trajectory.  MAC
        tracking provides the safeguard that prevents spurious modes from
        being misidentified as the tracked structural mode.
  \item Note that the sensitivity formula Eq.~(5) involves only the
        stiffness matrix derivative and the static displacement field,
        neither of which is dominated by low-density element contributions.
\end{itemize}

\newpage
% ============================================================
\section{Minor Revisions (C-group)}
% ============================================================

\subsection{C1 — Application references in Introduction}

Add two or three references to engineering applications of frequency
maximization in the first paragraph of the Introduction
(e.g.\ aerospace panels, seismic design of civil structures, rotating
machinery foundations). Suitable candidates from the existing literature
include Ferrari et al.~(2018), Lee et al.~(2024), and Wu and Li~(2024),
which are already cited later in Section~2.

% -----------------------------------------------------------
\subsection{C2 — SIMP, BESO, LSM, MMC references at first mention}

The reviewer requests canonical references when these methods are first
introduced. The following additions are needed (in Introduction, second
paragraph):

\begin{itemize}[leftmargin=1.6em]
  \item \textbf{SIMP:} Bends\o{}e and Sigmund (1999), \emph{Arch.\ Appl.\
        Mech.}
  \item \textbf{LSM:} Wang, Wang, and Guo (2003), \emph{Comput.\ Methods
        Appl.\ Mech.\ Eng.}; and Allaire, Jouve, and Toader (2004),
        \emph{J.\ Comput.\ Phys.}
  \item \textbf{ESO/BESO:} Xie and Steven (1993), \emph{Comput.\ Struct.};
        Querin et al.\ (2000), \emph{Comput.\ Methods Appl.\ Mech.\ Eng.}
  \item \textbf{MMC:} Guo et al.\ (2014), \emph{J.\ Appl.\ Mech.}
\end{itemize}

Add the corresponding BibTeX entries to \texttt{literature.bib}.

% -----------------------------------------------------------
\subsection{C3 — Rayleigh's principle reference and introductory note}

Add a reference for Rayleigh's principle at its first mention in the
Introduction (currently line~183 of \texttt{main.tex}). A suitable
general reference is Strang and Fix \emph{An Analysis of the Finite Element
Method} or any standard structural dynamics textbook. Add one or two sentences
explaining the principle for readers unfamiliar with it: the Rayleigh quotient
$R(\mathbf{u}) = \mathbf{u}^\top\mathbf{K}\mathbf{u}/(\mathbf{u}^\top\mathbf{M}\mathbf{u})$ provides an upper bound for the
fundamental eigenfrequency for any nonzero displacement field~$\mathbf{u}$.

% -----------------------------------------------------------
\subsection{C4 — MMA reference in Section~2.1}

The paper mentions MMA in Section~2.1 (Related Work, Olhoff subsection)
but does not add the Svanberg~(1987) reference at the first occurrence in
that section. Verify that the citation \texttt{\textbackslash citep\{Svanberg1987\}}
is present at every first mention of MMA in both Sections~2.1 and~3 (clamped
beam).

% -----------------------------------------------------------
\subsection{C5 — Hardware and platform details}

Add a footnote or a sentence at the start of Section~4 (Numerical examples /
Computational aspects) stating the computational platform used for all
benchmarks: CPU model, number of cores used (serial or parallel), available
RAM, operating system, MATLAB version (MATLAB 2025b is already mentioned
on line~597), and Python version if the Python code was also used to
generate Table~1 data.

\textbf{Audit note:} do not state Python/NumPy/SciPy versions unless the
table is actually regenerated with the Python implementation.  The manuscript
currently states that all results were generated and verified in MATLAB
2025b, and the active performance driver
\texttt{examples/Performance/performance\_comparison.m} calls the MATLAB JSON
wrapper.  The platform statement must match the implementation used for the
final Table~1.

% -----------------------------------------------------------
\subsection{C6 — Sensitivity filter equation in Section~2.1}

Reviewer~3 requests an equation and reference for the sensitivity filter
in Section~2.1 (Related Work) where filtering is first mentioned. The
filter equation is already present in Section~2.4 (Sensitivity filtering
and design update). Add a forward reference: ``\ldots a density-weighted
spatial filter of radius $r_{\min}$ (see Section~2.4 and
\cite{Sigmund2007}) is applied\ldots''. The reference Sigmund (2007)
\emph{Struct.\ Multidisc.\ Optim.} is a standard citation for the
sensitivity filter.

% -----------------------------------------------------------
\subsection{C7 — Explicit sensitivity for optimization problem (4)}

Reviewer~3 comment~\#7 requests the explicit sensitivity expression for
the compliance objective in Eq.~(4). This is already given as Eq.~(5) in
Section~2.2. A cross-reference at the end of the problem statement in
Eq.~(4) is sufficient: ``The sensitivity of $c$ with respect to $x_e$ is
derived in Eq.~(5) below.''

% -----------------------------------------------------------
\subsection{C8 — Supplementary material clarification}

Reviewer~3 comment~\#13 asks what the supplementary material contains.
Add a sentence in the Data Availability or a footnote clarifying:
``Supplementary material includes additional convergence plots and MAC
correlation matrices for all benchmark cases; it is available in the
public repository.''  If no formal supplementary file was submitted,
remove or correct any implicit reference to it.

% -----------------------------------------------------------
\subsection{C9 — Load-sensitivity treatment (Reviewer~4, comment~\#1)}

Reviewer~4 asks about lines 798--799 of \texttt{main.tex} (the sentence
concerning sensitivity computation of the right-hand side of the
equivalent static problem). The paper already states that
$\partial\mathbf{f}/\partial x_e$ is \emph{omitted} from the total
derivative of $c$ (Section~2.2, after Eq.~4). Add the explicit statement
``$\partial\mathbf{f}/\partial x_e$ is \textbf{not} computed in the
proposed method'' at the specific location the reviewer references, and
add a sentence justifying this omission (the frozen load is treated as
externally specified, so its density gradient is set to zero, yielding the
standard stiffness-based compliance sensitivity Eq.~(5)).

\textbf{Code reference:} in both MATLAB and Python solvers, the
\texttt{semi\_harmonic} branch returns zero load-sensitivity contribution.
The \texttt{harmonic} branch is different: it computes a
$\partial\mathbf{M}/\partial x_e$ load term while ignoring
$\partial\omega/\partial x_e$ and
$\partial\boldsymbol{\Phi}/\partial x_e$.  Therefore the paper's statement
that $\partial\mathbf{f}/\partial x_e$ is omitted is correct only for the
\texttt{semi\_harmonic} path used in the reported examples.  If a refresh
study uses \texttt{harmonic}, add a switch to zero its load-sensitivity term
or state explicitly that the refresh study changes both the update schedule
and the sensitivity model.

% -----------------------------------------------------------
\subsection{C10 — Numerical difficulties from design-dependent load (Reviewer~4, comment~\#2)}

Reviewer~4 asks whether the design-dependent inertial load introduces
numerical difficulties analogous to self-weight optimization.

Add a short remark in Section~2.2 or Section~4.2: after the load definition
is corrected in A0, in the proposed frozen-mode method the load
$\mathbf{f}(\mathbf{x})$ is design-dependent only through the density-scaled
mass/load-amplitude term. Because
$\boldsymbol{\Phi}_1$ is fixed, the load varies linearly with density at
each element DOF, analogous to a gravitational body force. The load sensitivity
is zeroed (see C9), so no additional numerical coupling between the load and
the stiffness update arises. No convergence difficulties attributable to this
coupling were observed in the numerical examples. In contrast, the Yuksel
iterative scheme \emph{does} couple the load estimate and the design update at
each iteration, which Munro~(2024) identifies as the cause of the additional
bisection correction step needed in self-weight OC implementations.

% -----------------------------------------------------------
\subsection{C11 — Applicability to symmetric vs.\ asymmetric domains (Reviewer~4, comment~\#4)}

Reviewer~4 asks whether the method is better suited to domains with symmetry
(where repeated eigenvalues arise) or to asymmetric domains (where a dominant
mode exists naturally).

Add a remark in the Conclusions or Discussion: the method targets a
\emph{specific} mode shape (frozen from the declared reference design), so it
does not require the mode to be dominant. However, for symmetric domains, the
initial eigenproblem may produce degenerate eigenvalues, making the chosen
eigenvector non-unique. The multi-mode weighting with parameter~$\alpha$
already handles the two-mode case; for fully symmetric domains with more
degenerate modes, the choice of $\boldsymbol{\Phi}_j$ may be arbitrary and
the load should be constructed as a superposition of the degenerate modes.

\newpage
% ============================================================
\section{Action Plan and Timeline}
% ============================================================

\begin{table}[htbp]
\centering
\caption{Proposed revision timeline.}
\label{tab:timeline}
\renewcommand\arraystretch{1.2}
\begin{tabularx}{\textwidth}{@{}p{0.5cm}p{2.5cm}p{9cm}p{1.8cm}@{}}
\toprule
ID & Owner & Task & Target \\
\midrule
A0 & Tauzowski & Resolve load-definition and baseline mismatch; rerun affected results if code changes & Week 1 \\
A1 & Tauzowski & Rewrite abstract + Section~2.4 last paragraphs & Week 1 \\
A2 & Tauzowski & Add init-time column to Table~1; rerun benchmarks with timing & Week 1 \\
A3 & Szczepańczyk & Diagnose Yuksel iteration discrepancy; add explanation & Week 1 \\
A4 & Błachowski  & Implement controlled N-refresh study with matched sensitivity model & Week 2 \\
A5 & Tauzowski & Add mode-coalescence and a~posteriori check discussion & Week 1 \\
A6 & Tauzowski & Revise efficiency claim language in abstract and Section~4 & Week 1 \\
\midrule
B1 & Szczepańczyk & Visualise and classify two new modes at $\alpha=1$ & Week 2 \\
B2 & Szczepańczyk & Add note on high-index MAC tracking + spurious modes & Week 2 \\
B3 & Szczepańczyk & Apply B1/B2 discussion to building example & Week 2 \\
B4 & Tauzowski & Add RAMP discussion; optionally run post-proc. RAMP eigensolve & Week 3 \\
B5 & Błachowski  & Run Heaviside version for one case; add note on grey regions & Week 2 \\
B6 & Tauzowski & Add explanation of spurious-mode robustness to Section~2.2 & Week 1 \\
\midrule
C1--C8 & Błachowski  & Reference additions, filter equation, platform details & Week 1 \\
C9     & Tauzowski   & Clarify load-sensitivity treatment; check code vs.\ paper & Week 1 \\
C10    & Tauzowski   & Add remark on design-dependent load difficulty & Week 1 \\
C11    & Błachowski  & Add note on symmetric vs.\ asymmetric domains & Week 1 \\
\midrule
All    & All         & Internal review of revised manuscript & Week 3 \\
All    & All         & Prepare response letter and resubmit & Week 4 \\
\bottomrule
\end{tabularx}
\end{table}

% ============================================================
\section{Code Changes Required}
% ============================================================

The following source code changes are needed to support the revised
experiments requested by the reviewers.

\begin{enumerate}[leftmargin=1.6em]

  \item \textbf{Load-definition alignment}
        (\texttt{analysis/ourApproach/Matlab/topopt\_freq.m},
        \texttt{analysis/ourApproach/Python/topopt\_freq.py}, and example
        JSON files).  Resolve A0 before running any new experiments.  Either
        change \texttt{semi\_harmonic} to use
        $(\omega_j^{(0)})^2\mathbf{M}_0\boldsymbol{\Phi}_j$ and rerun
        reported results, or revise the manuscript to match the current
        $\omega_j^{(0)}\mathbf{M}_0\boldsymbol{\Phi}_j$ reference-load
        implementation.  Also decide whether the reference baseline is
        \texttt{"solid"} or \texttt{"initial"} and make all examples and text
        consistent.

  \item \textbf{Timing instrumentation}
        (MATLAB and Python solvers, plus JSON wrappers if used for Table~1).
        Wrap the active initialization eigensolves with timers and return
        \texttt{init\_time}, \texttt{loop\_time}, and \texttt{total\_time} in
        diagnostics.  Disable \texttt{save\_frequency\_iterations} during
        benchmark timing runs, because it forces per-iteration eigensolves in
        the proposed MATLAB solver.

  \item \textbf{N-refresh benchmark script}.
        Create a controlled refresh-study script for the simply-supported
        beam at $400{\times}50$.  Do not rely directly on the existing
        \texttt{harmonic/update\_after} path unless a switch is added to keep
        the same load formula and zero-load-sensitivity approximation as the
        proposed method.  Record \texttt{n\_iter}, \texttt{t\_iter},
        \texttt{init\_time}, total time, and final $\omega_1$ in a CSV file.

  \item \textbf{Mode-shape export for B1/B2}.
        The MATLAB wrapper already contains postprocessing routines for
        reference/topology modes, correlation matrices, and mode plots.  Use
        those first.  Only extend the Python \texttt{info} dictionary if the
        revised experiments will be regenerated through the Python path.

  \item \textbf{RAMP post-processing eigensolve (optional, for B4)}.
        Implement only after selecting the exact RAMP stiffness/mass
        interpolation pair from the cited literature.  Do not add an ad hoc
        \texttt{ramp\_q} formula without validating that it actually shifts
        low-density localized modes upward for this model.

\end{enumerate}

% ============================================================
\section{New Experiments Required}
% ============================================================

\begin{enumerate}[leftmargin=1.6em]

  \item \textbf{Eigenpair refresh sensitivity study} (required by A4).
        Simply-supported beam, $400{\times}50$, five values of
        refresh interval: $\infty$ (frozen), 50, 10, 5, and 1 (full
        true-eigenpair refresh).  Report final
        $\omega_1$, runtime, and iteration count.  Expected effort:
        $\sim$10 runs, each $<$5~min.

  \item \textbf{Mode-shape visualisation for the two new modes at
        $\alpha=1$} (required by B1).  No new optimization run needed;
        reload the existing density field and call \texttt{eigsh} for
        the first 5 modes. Plot mode shapes of modes~1 and~2 and
        include in supplementary material or in the paper.

  \item \textbf{Heaviside-projection run for one clamped-beam case}
        (optional for B5, $\alpha=1$ only).  Set
        \texttt{"optimization.filter.heaviside": true} and
        \texttt{"heaviside\_beta\_schedule": [1,2,4,8,16,32,64]}
        in \texttt{BeamTopOptFreq.json}; if using MATLAB, first forward this
        parsed flag into the \texttt{ourApproach} run configuration.
        Expected effort: one run, $\sim$10~min.

\end{enumerate}

\newpage
% ============================================================
\section{Algorithm Pseudocodes and Complexity Analysis}
\label{sec:algorithms}
% ============================================================

This section presents pseudocodes for the three compared methods in a
uniform notation, followed by a detailed complexity assessment.
The comparison must distinguish canonical algorithms from the repository
implementations used to generate Table~1.  In particular, the local
Olhoff implementation adds a trial eigensolve after each MMA update, and the
local Yuksel timing aggregates Stage~1 and Stage~2 iterations.  Those
implementation choices should be stated explicitly wherever empirical
runtime ratios are discussed.
Notation throughout: $N_e$ = number of finite elements;
$N_\mathrm{dof} \approx 2N_e$ = degrees of freedom;
$J$ = number of candidate eigenpairs tracked;
$p,d$ = SIMP penalization exponents for stiffness and mass;
$\varepsilon$ = convergence tolerance;
$bw \approx \sqrt{N_e}$ = effective bandwidth of the stiffness matrix
for a 2-D quasi-banded mesh.

% -----------------------------------------------------------
\subsection{Algorithm 1: Olhoff--Du Bound Formulation}
\label{subsec:alg-olhoff}
% -----------------------------------------------------------

The method maximizes a scalar bound variable $\beta \leq \omega_j^2$,
$j = n,\ldots,J$, through a nested double-loop procedure.
The outer loop solves the eigenvalue problem and assembles sensitivities;
the inner loop runs MMA iteratively until the design increments converge
\cite{Olhoff2014}.

\medskip
\begin{algobox}
\kw{function} \fn{Olhoff-Du}(\var{mesh}, \var{J}, \var{$V_f$})
\kw{returns} optimal density field $x^*$\\[2pt]
\ind\kw{inputs}: \var{mesh}, finite element mesh and boundary conditions\\
\ind\hspace*{4.2em}\var{J}, number of candidate eigenfrequencies ($J \geq n$)\\
\ind\hspace*{4.2em}$V_f$, prescribed volume fraction\\[2pt]
\ind\kw{local variables}:
  $x = \{x_e\}$, design variables, initially $x_e = V_f$\\
\ind\hspace*{9.5em}
  $\beta$, scalar bound variable\\
\ind\hspace*{9.5em}
  $\lambda_j, \omega_j, \boldsymbol{\varphi}_j$,
  eigenvalues, frequencies, eigenvectors, $j=1,\ldots,J$\\
\ind\hspace*{9.5em}
  $N$, detected multiplicity of $\omega_n$\\
\ind\hspace*{9.5em}
  $\mathbf{f}_{sk}$, generalized gradient vectors (or
  $\nabla\!\lambda_j$ if $N=1$)\\
\ind\hspace*{9.5em}
  $\Delta\var{x}$, vector of design variable increments\\[6pt]
%
$x_e \assign V_f$ \quad for all $e = 1,\ldots,N_e$;
\quad $\beta \assign 0$\\[4pt]
%
\kw{repeat} \cm{// outer (main) loop}\\
\ind\cm{// Step 1 — Eigensolve}\\
\ind solve $\mathbf{K}(\var{x})\,\boldsymbol{\varphi}_j
           = \omega_j^2\,\mathbf{M}(\var{x})\,\boldsymbol{\varphi}_j$,
      $\boldsymbol{\varphi}_j^{\top}\mathbf{M}\boldsymbol{\varphi}_k=\delta_{jk}$,
      $j=1,\ldots,J$\\
\ind detect multiplicity $N$ of $\omega_n$ (and $R$ of $\omega_{n-1}$)\\[3pt]
%
\ind\cm{// Step 2 — Sensitivity analysis}\\
\ind\kw{if} $N > 1$ \kw{then}\\
\indd compute generalized gradients
      $\mathbf{f}_{sk}=\{\boldsymbol{\varphi}_s^{\top}
      (\mathbf{K}'_{\rho_e}-\tilde{\lambda}\mathbf{M}'_{\rho_e})
      \boldsymbol{\varphi}_k\}^{\top}$,\;
      $s,k=n,\ldots,n{+}N{-}1$\\
\ind\kw{else}\\
\indd compute standard gradients
      $(\lambda_j)'_{\rho_e}=
      \boldsymbol{\varphi}_j^{\top}
      (\mathbf{K}'_{\rho_e}-\lambda_j\mathbf{M}'_{\rho_e})
      \boldsymbol{\varphi}_j$,\;
      $j=n,\ldots,J$\\[3pt]
%
\ind\cm{// Step 3 — Inner loop: MMA sub-problem for design increments}\\
\ind\kw{repeat} \cm{// inner loop}\\
\indd solve MMA sub-problem (bound form, Eq.~19 of \cite{Olhoff2014})
      for increments $\Delta x_e$,\; $e=1,\ldots,N_e$\\
\ind\kw{until} $\|\Delta\var{x}\| < \varepsilon_{\mathrm{inner}}$\\[3pt]
%
\ind\cm{// Step 4 — Update design variables}\\
\ind $x_e \assign x_e + \Delta x_e$ \quad for all $e$\\
\ind $\beta \assign \min_{j=n,\ldots,J}\omega_j^2(\var{x})$\\[4pt]
%
\kw{until} $\|\Delta\var{x}\| < \varepsilon$\\[4pt]
\kw{return} $x$
\end{algobox}

% -----------------------------------------------------------
\subsection{Algorithm 2: Yuksel--Yilmaz Two-Stage Static Method}
\label{subsec:alg-yuksel}
% -----------------------------------------------------------

Two sequential static compliance minimizations replace the per-iteration
eigensolve.
Stage~1 produces a static displacement field as the initial eigenvector
estimate; Stage~2 iteratively refines this estimate through a
design-dependent inertial load \cite{Yuksel2025}.

\medskip
\begin{algobox}
\kw{function} \fn{Yuksel-Yilmaz}(\var{mesh}, $V_f$,
$\mathbf{f}_s$) \kw{returns} optimal density field $x^*$\\[2pt]
\ind\kw{inputs}: \var{mesh}, finite element mesh and boundary conditions\\
\ind\hspace*{4.2em}$V_f$, prescribed volume fraction\\
\ind\hspace*{4.2em}$\mathbf{f}_s$, static point load
  (midpoint for symmetric BCs; from solid eigensolve for asymmetric BCs)\\[2pt]
\ind\kw{local variables}:
  $x=\{x_e\}$, design variables, initially $x_e = V_f$\\
\ind\hspace*{9.5em}
  $\mathbf{u}_s$, static displacement from Stage~1\\
\ind\hspace*{9.5em}
  $\mathbf{u}_d$, current mode-shape estimate;
  $\hat{\mathbf{u}}_d$, its unit vector\\
\ind\hspace*{9.5em}
  $\mathbf{f}_d$, design-dependent inertial load vector\\[6pt]
%
$x_e \assign V_f$ \quad for all $e$\\[4pt]
%
\cm{// Stage 1 — mode-shape estimation via static compliance}\\
\kw{repeat}\\
\ind solve $\mathbf{K}(\var{x})\,\mathbf{u}_s = \mathbf{f}_s$\\
\ind compute
    $\frac{\partial c}{\partial x_e}
    = p x_e^{p-1}(E_0-E_{\min})\,\mathbf{u}_{s,e}^{\top}\mathbf{K}_e^0\,\mathbf{u}_{s,e}$,
    \quad $e=1,\ldots,N_e$\\
\ind apply sensitivity filter; update $x$ by OC bisection on
    $V(\var{x})=V_f\cdot V_0$\\
\kw{until} $\max_e|\Delta x_e| < \varepsilon_1$\\[4pt]
%
$\mathbf{u}_d \assign \mathbf{u}_s$
\quad\cm{// initialise eigenvector estimate from Stage~1 result}\\[4pt]
%
\cm{// Stage 2 — self-consistent iterative compliance with inertial load}\\
\kw{repeat}\\
\ind $\hat{\mathbf{u}}_d \assign \mathbf{u}_d\,/\,\|\mathbf{u}_d\|$
  \quad\cm{// normalise current mode estimate}\\
\ind $\mathbf{f}_d \assign \mathbf{M}(\var{x})\,\hat{\mathbf{u}}_d$
  \quad\cm{// assemble design-dependent inertial load}\\
\ind solve $\mathbf{K}(\var{x})\,\mathbf{u}_d = \mathbf{f}_d$\\
\ind compute $\frac{\partial c}{\partial x_e}
    = p x_e^{p-1}(E_0-E_{\min})\,\mathbf{u}_{d,e}^{\top}\mathbf{K}_e^0\,\mathbf{u}_{d,e}$\\
\ind apply sensitivity filter; update $x$ by OC bisection on
    $V(\var{x})=V_f\cdot V_0$\\
\kw{until} $\max_e|\Delta x_e| < \varepsilon_2$\\[4pt]
%
apply Heaviside post-processing to obtain solid--void design\\
solve $\mathbf{K}(\var{x}^*)\boldsymbol{\varphi}
      = \omega_1^2\,\mathbf{M}(\var{x}^*)\boldsymbol{\varphi}$
\quad\cm{// final eigensolve to report $\omega_1$}\\[4pt]
\kw{return} $x^*$
\end{algobox}

% -----------------------------------------------------------
\subsection{Algorithm 3: Proposed Quasi-Static Approximation}
\label{subsec:alg-proposed}
% -----------------------------------------------------------

A single eigensolve on the declared reference domain freezes the eigenpair
permanently.  The pseudocode below assumes the physical
$(\omega_j^{(0)})^2\mathbf{M}\boldsymbol{\Phi}_j$ inertial-load definition;
if A0 is resolved by retaining the current \texttt{semi\_harmonic}
implementation, replace this formula consistently with the implemented
reference-load definition.
Every subsequent iteration is a standard compliance minimization step;
the multi-mode load-weighting parameter~$\alpha$ provides simultaneous
two-mode control without any structural change to the algorithm.

\medskip
\begin{algobox}
\kw{function} \fn{Quasi-Static-Freq}(\var{mesh}, $V_f$,
$n_{\mathrm{modes}}$, $\alpha$) \kw{returns} optimal density field $x^*$\\[2pt]
\ind\kw{inputs}: \var{mesh}, finite element mesh and boundary conditions\\
\ind\hspace*{4.2em}$V_f$, prescribed volume fraction\\
\ind\hspace*{4.2em}$n_{\mathrm{modes}}$, number of modes to target
  (1 or 2)\\
\ind\hspace*{4.2em}$\alpha \in [0,1]$, load-weighting parameter\\[2pt]
\ind\kw{local variables}:
  $x=\{x_e\}$, design variables, initially $x_e=V_f$\\
\ind\hspace*{9.5em}
  $(\omega_j^{(0)},\boldsymbol{\Phi}_j)$,
  reference eigenpairs $j=1,\ldots,n_{\mathrm{modes}}$,
  \emph{fixed throughout}\\
\ind\hspace*{9.5em}
  $\mathbf{f}_\alpha$, weighted inertial load vector\\
\ind\hspace*{9.5em}
  $\mathbf{u}$, static displacement field\\[6pt]
%
\cm{// Initialisation — single eigensolve on declared reference domain}\\
$x_e \assign V_f$ \quad for all $e = 1,\ldots,N_e$\\
assemble $\mathbf{K}_0$ and $\mathbf{M}_0$ at the chosen reference design\\
solve $\mathbf{K}_0\boldsymbol{\Phi}_j
      = (\omega_j^{(0)})^2\mathbf{M}_0\boldsymbol{\Phi}_j$
  \quad for $j=1,\ldots,n_{\mathrm{modes}}$
  \quad\cm{// performed \emph{once}, then $(\omega_j^{(0)},\boldsymbol{\Phi}_j)$ frozen}\\[4pt]
%
\cm{// Main loop — pure static compliance minimization, no eigensolve}\\
\kw{repeat}\\
\ind assemble $\mathbf{K}(\var{x})$ and $\mathbf{M}(\var{x})$\\
\ind $\mathbf{f}_\alpha \assign
  \alpha\,(\omega_1^{(0)})^2\mathbf{M}(\var{x})\boldsymbol{\Phi}_1
  + (1{-}\alpha)\,(\omega_2^{(0)})^2\mathbf{M}(\var{x})\boldsymbol{\Phi}_2$
  \quad\cm{// frozen $\boldsymbol{\Phi}_j$}\\
\ind solve $\mathbf{K}(\var{x})\,\mathbf{u} = \mathbf{f}_\alpha(\var{x})$\\
\ind compute
  $\dfrac{\partial c}{\partial x_e}
  = -p\,x_e^{p-1}(E_0-E_{\min})\,\mathbf{u}_e^{\top}\mathbf{K}_e^0\,\mathbf{u}_e$
  \quad\cm{// load sensitivity $\partial\mathbf{f}/\partial x_e$ set to zero}\\
\ind apply density-weighted sensitivity filter of radius $r_{\min}$\\
\ind update $x$ by OC bisection on volume constraint
  $\sum_e x_e V_e = V^*$\\
\kw{until} $\max_e|\Delta x_e| < \varepsilon_{\mathrm{tol}}$\\[4pt]
%
\cm{// Post-analysis}\\
solve $\mathbf{K}(\var{x}^*)\boldsymbol{\varphi}_j
      = \omega_j^2\,\mathbf{M}(\var{x}^*)\boldsymbol{\varphi}_j$,
  \;$j=1,\ldots,3$
  \quad\cm{// one final eigensolve to verify $\omega_1$}\\[4pt]
\kw{return} $x^*$
\end{algobox}

% -----------------------------------------------------------
\subsection{Complexity Analysis and Comparison}
\label{subsec:complexity}
% -----------------------------------------------------------

\subsubsection*{Dominant operations}

For a 2-D FEM mesh with $N_e$ elements the key cost drivers per
iteration are listed below.
Sparse Cholesky factorization of a quasi-banded system scales as
$O(N_e^{1.5})$ and dominates all other operations for meshes beyond
a few thousand elements.
An eigensolve via shift-and-invert Arnoldi
(\textsc{arpack}/\texttt{eigsh}) requires one Cholesky factorization
plus $k_{\mathrm{Krylov}} \approx 3J$ Krylov iterations, each costing
$O(N_e \cdot bw) = O(N_e^{1.5})$ for a banded system; the
constant~$C_{\mathrm{eig}}$ therefore exceeds the simple-solve
constant~$C_{\mathrm{solve}}$ by a factor of $3J$--$5J$ in practice.

\begin{table}[htbp]
\centering
\caption{Per-iteration operation count.
$bw \approx \sqrt{N_e}$ denotes the effective bandwidth;
$J$ is the number of tracked modes;
$c = r_{\min}/h$ is the filter radius in elements (constant).}
\label{tab:ops}
\renewcommand\arraystretch{1.2}
\begin{tabularx}{\textwidth}{@{}lXXX@{}}
\toprule
Operation & Olhoff--Du implementation & Yuksel--Yilmaz (Stage~2) & Proposed \\
\midrule
Shift-invert eigensolves / factorizations
  & $2 \times O(N_e^{1.5})$ (main + trial)
  & $0$ & $0$ \\
Static linear solves
  & --- & $1 \times O(N_e^{1.5})$ & $1 \times O(N_e^{1.5})$ \\
Stiffness assembly & $O(N_e)$ & $O(N_e)$ & $O(N_e)$ \\
Mass assembly & $O(N_e)$ & $O(N_e)$ & $O(N_e)$ \\
Sensitivity ($J$ modes) & $J \times O(N_e)$ & $O(N_e)$ & $O(N_e)$ \\
Filter (conv.) & $J \times O(c^2 N_e)$ & $O(c^2 N_e)$ & $O(c^2 N_e)$ \\
MMA update & $O((J{+}2)N_e)$ & --- & --- \\
OC bisection & --- & $O(N_e)$ & $O(N_e)$ \\
\midrule
\textbf{Total} & $\boldsymbol{O(N_e^{1.5})}$ (large const.)
               & $\boldsymbol{O(N_e^{1.5})}$ (small const.)
               & $\boldsymbol{O(N_e^{1.5})}$ (small const.) \\
\bottomrule
\end{tabularx}
\end{table}

\subsubsection*{Total cost and empirical validation}

The total wall-clock time is $T = n_{\mathrm{iter}} \times
t_{\mathrm{iter}}$, where $t_{\mathrm{iter}}$ is the per-iteration
time and $n_{\mathrm{iter}}$ is the number of iterations to
convergence.
Table~\ref{tab:complexity} records both quantities from the benchmark
at four mesh resolutions (Table~1 of the manuscript) and derives
the implied speedup ratios.

\begin{table}[htbp]
\centering
\caption{Empirical performance at the four mesh sizes.
$t_{\mathrm{iter}}$ = average time per iteration (s);
$n_{\mathrm{iter}}$ = total iterations to convergence;
$T$ = total wall-clock time (s);
speedup is relative to Olhoff--Du.
Data reproduced from Table~1 of the manuscript; final values should be
updated after rerunning with timing diagnostics separated and
\texttt{save\_frequency\_iterations=false}.}
\label{tab:complexity}
\renewcommand\arraystretch{1.2}
\setlength{\tabcolsep}{4pt}
\begin{tabular}{@{}l r r r r r r r r r@{}}
\toprule
& \multicolumn{3}{c}{\textbf{Olhoff--Du}}
& \multicolumn{3}{c}{\textbf{Yuksel--Yilmaz}}
& \multicolumn{3}{c}{\textbf{Proposed}} \\
\cmidrule(lr){2-4}\cmidrule(lr){5-7}\cmidrule(lr){8-10}
Mesh & $t_\mathrm{iter}$ & $n_\mathrm{iter}$ & $T$ (s)
     & $t_\mathrm{iter}$ & $n_\mathrm{iter}$ & $T$ (s)
     & $t_\mathrm{iter}$ & $n_\mathrm{iter}$ & $T$ (s) \\
\midrule
$160{\times}20$  & 0.09 & 372 & 32.0
                 & 0.02 & 300 &  4.7
                 & 0.02 & 264 &  5.1 \\
$240{\times}30$  & 0.20 & 363 & 74.4
                 & 0.03 & 702 & 21.5
                 & 0.04 & 250 & 10.5 \\
$320{\times}40$  & 0.39 & 350 & 135.7
                 & 0.05 & 1723 & 89.4
                 & 0.08 & 217 & 16.4 \\
$400{\times}50$  & 0.64 & 345 & 221.9
                 & 0.08 & 1083 & 84.1
                 & 0.12 & 263 & 31.1 \\
\midrule
Scaling exponent
  & \multicolumn{3}{c}{$t_\mathrm{iter}\sim O(N_e^{1.30})$}
  & \multicolumn{3}{c}{$t_\mathrm{iter}\sim O(N_e^{1.25})$}
  & \multicolumn{3}{c}{$t_\mathrm{iter}\sim O(N_e^{1.35})$} \\
\bottomrule
\end{tabular}
\end{table}

The per-iteration scaling exponent is estimated by a log--log fit
across the four mesh sizes.
All three methods exhibit an empirical exponent near~$1.3$, consistent
with the theoretical $O(N_e^{1.5})$ bound for banded Cholesky (the
incomplete fill-in of the actual sparse pattern yields a lower
empirical exponent).

\subsubsection*{Structured complexity summary}

Table~\ref{tab:summary} provides a qualitative side-by-side comparison
of the three algorithms covering algorithmic structure, achieved
accuracy, and practical properties.

\begin{table}[htbp]
\centering
\caption{Qualitative algorithm comparison.
``$\approx$'' denotes approximate (heuristic) convergence criterion.}
\label{tab:summary}
\renewcommand\arraystretch{1.3}
\begin{tabularx}{\textwidth}{@{}p{4.5cm}XXX@{}}
\toprule
Property & Olhoff--Du & Yuksel--Yilmaz & Proposed \\
\midrule
\textbf{Eigensolves in loop} & 2 per outer iteration
  & 0 (Stage~2) & 0 \\
\textbf{One-time eigensolves} & 0--1 (diagnostic)
  & 0--1 (asymm.\ BCs) & 1 (init, fixed) \\
\textbf{Linear solves / iter.} & 2 (main + trial)
  & 1 & 1 \\
\textbf{Inner optimization loop} & Yes (MMA, $\sim$5--20 iters)
  & No & No \\
\textbf{Mode shape updated} & Every iteration (new eigenvec.)
  & Every iteration (self-consistent) & Never (frozen) \\
\textbf{Handles repeated $\omega_j$} & Yes (sub-eigenproblem)
  & Partial (OC heuristic) & N/A (no eigensolve) \\
\textbf{Multi-mode targeting} & Yes ($J$ modes, bound $\beta$)
  & No (mode~1 only) & Yes ($\alpha$-weighted) \\
\textbf{Per-iter cost (observed)} & $\sim8\times$ baseline
  & $\sim1\times$ baseline & $\sim1.5\times$ baseline \\
\textbf{Typical iterations ($400{\times}50$)} & $\sim$345
  & $\sim$1083 (Stage~1+2) & $\sim$263 \\
\textbf{Total time speedup vs.\ Olhoff} & 1$\times$ (reference)
  & $2.6\times$ & $7.1\times$ \\
\textbf{Source of speedup} & ---
  & Low $t_\mathrm{iter}$, many iters & Moderate $t_\mathrm{iter}$, few iters \\
\textbf{Peak RAM ($400{\times}50$)} & 303~MB
  & 131~MB & 56~MB \\
\textbf{Projection / post-proc.} & Yes (Heaviside $\beta$-schedule)
  & Yes (Heaviside post-proc.) & Optional \\
\textbf{Achievable $\omega_1$ quality} & Near-optimal (bound)
  & Approx.\ (${\sim}8\%$ below bound) & Approx.\ (${\sim}9\%$ below bound) \\
\textbf{Code complexity} & High (MMA, Heaviside, bound var.)
  & Medium (two-stage logic) & Low (compliance + OC) \\
\bottomrule
\end{tabularx}
\end{table}

\subsubsection*{Key observations}

\begin{enumerate}[leftmargin=1.6em]

  \item \textbf{All three methods have the same asymptotic per-iteration
        complexity} $O(N_e^{1.5})$.
        The practical speedup of the Proposed method over Olhoff--Du
        ($7.1\times$ at $400{\times}50$) is therefore entirely determined
        by two factors: a smaller constant in the per-iteration cost
        (no Arnoldi iterations, no bound/MMA overhead, no trial eigensolve)
        and a substantially
        lower iteration count.

  \item \textbf{The dominant cost driver in Olhoff--Du is the double
        eigensolve per outer iteration.}
        At $400{\times}50$, each eigensolve requires a sparse Cholesky
        factorization of a $40{,}902 \times 40{,}902$ system, repeated
        twice (main + trial), accounting for $\sim$80\,\% of the
        $0.64$~s per-iteration time.
        The static-load methods replace both eigensolves in this implementation
        with a single linear solve, reducing the dominant cost by a
        factor of $\sim$5--8$\times$.

  \item \textbf{The Proposed method converges in fewer iterations than
        Yuksel--Yilmaz.}
        The fixed frozen gradient provides a consistent, non-oscillating
        descent direction from the first iteration, whereas the
        self-consistent load update in Yuksel--Yilmaz can introduce
        oscillatory behaviour that inflates the iteration count at fine
        meshes (1083 vs.\ 263 at $400{\times}50$, a ratio of $4.1\times$).
        The net speedup over Yuksel--Yilmaz is $84.1/31.1 = 2.7\times$,
        which is the product of the per-iteration cost ratio
        ($0.08/0.12 \approx 0.67$, i.e.\ Yuksel is slightly cheaper per
        iteration) and the iteration-count ratio ($1083/263 \approx 4.1$).
        The smallest mesh is an exception in the current table: Yuksel is
        slightly faster than the proposed method at $160{\times}20$.

  \item \textbf{Memory scales as $O(N_e)$ for all methods, but with
        very different constants.}
        The Olhoff--Du implementation retains two Cholesky-factorized
        matrices simultaneously (current + trial), MMA history arrays
        ($\mathbf{x}^{(k-1)},\mathbf{x}^{(k-2)}$, asymptote vectors),
        and the $(J{+}2)\times N_e$ gradient matrix, together accounting
        for the observed 303~MB at $400{\times}50$.
        The Proposed method needs only one factorized $\mathbf{K}_f$, the
        frozen vectors $\boldsymbol{\Phi}_j$ ($n_{\mathrm{modes}} \times N_e$),
        and the OC bisection buffers, yielding 56~MB at the same mesh ---
        a $5.4\times$ reduction.

  \item \textbf{Accuracy of the static-approximation methods is
        fundamentally limited.}
        Because neither Yuksel--Yilmaz nor the Proposed method tracks the
        true eigenvalue objective at each iteration, they should not be
        claimed to recover the bound-formulation optimum in general.
        The gap ($\sim$8--9\,\% below the Olhoff--Du $\omega_1$ for the
        simply supported beam) is the observed price of the static
        approximation for this benchmark.  Reducing that gap may require
        eigenpair refreshes, mode aggregation, or returning to an eigenvalue
        objective, all of which change the cost/accuracy trade-off.

\end{enumerate}

\bibliographystyle{plain}
\bibliography{../literature}

\end{document}
